
Equivariant weight-space encoders exploit permutation symmetry in neural network weights by construction, yet their sample efficiency advantage over plain encoders has not been measured in a controlled, shared-split study. This work compares three conditions on the ModelZooDataset CIFAR-10 CNN zoo --- a flat multilayer perceptron (flat-MLP, plain baseline), a flat-MLP with permutation augmentation (PermAug), and GNN-NFN (structural equivariance) --- across training set sizes {100, 250, 500, 1000, full}. GNN-NFN reaches 90\% of its peak $R^2$ at approximately $N\approx 147$ training models; flat-MLP does not reach 90\% of its own peak $R^{2}=0.886$ within any sampled training size up to N=1,000, requiring the full dataset (~7,000 models), yielding an efficiency ratio of approximately $47\times$. This advantage is data-regime dependent: at N=100, PermAug ($R^{2}=0.138)$ outperforms structural equivariance (GNN-NFN $R^{2}=-0.016)$ --- a single-seed observation requiring multi-seed replication --- revealing a minimum-data threshold of approximately 150--250 models below which data augmentation of a plain encoder is superior. At full training scale, all encoders converge within 1\% $R^2$, consistent with recently proposed expressivity equivalence theory \citep{Dayan2026Expressive}. Permutation equivariance is verified to floating-point precision for GNN-NFN (max\_diff=1.$80\times 10^{-6}$, 10,000 checks) and DWSNets (max\_diff=7.$45\times 10^{-9}$, 1,000 checks). These findings provide a concrete decision boundary for practitioners choosing between encoder types under zoo-size constraints.

